\section{Types}
\label{sec:types}

The source syntax of types is given in \cref{fig:syntax}. This section assigns static
meaning to that syntax. The type system has three layers. First, contexts and
session types are quotiented by the equations in \cref{fig:syntax-types}, and
weakening is controlled by a subcontext preorder. Second, the predicates
$\Unr$, $\Mobile$, $\Bounded$, $\Skips$, and $\SNew$ classify values, contexts,
and sessions that may be duplicated, transported, split, or freshly created.
Third, the typing judgments for constants, expressions, and processes use
these classifications to enforce protocol fidelity and ownership transfer.

Typing contexts combine assumptions with two constructors. The ordered product
$\Ctx[1]\CtxSeq\Ctx[2]$ records that the resources in $\Ctx[1]$ must be used
before the resources in $\Ctx[2]$. The unordered product
$\Ctx[1]\CtxPar\Ctx[2]$ records independent resources. When a rule is
parametric in either constructor, we write $\ctxcolor{\Ctx[1]\CtxPatOp\Ctx[2]}$
for either $\Ctx[1]\CtxSeq\Ctx[2]$ or $\Ctx[1]\CtxPar\Ctx[2]$.

\begin{figure}
  \medskip
  Equality of typing contexts is the equivalence closure of the following
  axioms.
  \begin{mathpar}
    \CtxEmpty\CtxPar\Ctx = \Ctx \and
    \Ctx[1] \CtxPar \Ctx[2] = \Ctx[2] \CtxPar \Ctx[1] \and
    \Ctx[1] \CtxSeq (\Ctx[2] \CtxSeq \Ctx[3]) = (\Ctx[1] \CtxSeq \Ctx[2]) \CtxSeq \Ctx[3] \and
    \Ctx[1] \CtxPar (\Ctx[2] \CtxPar \Ctx[3]) = (\Ctx[1] \CtxPar \Ctx[2]) \CtxPar \Ctx[3] \\
    \inferrule{\Unr\Ctx}{\Ctx \CtxPar \Ctx = \Ctx} \and
    \inferrule{\Mobile\Ctx[1]}{\Ctx[1] \CtxPar \Ctx[2]
      = \Ctx[1] \CtxSeq \Ctx[2]} \and
    \inferrule{\Mobile\Ctx[2]}{\Ctx[1] \CtxPar \Ctx[2]
      = \Ctx[1] \CtxSeq \Ctx[2]}
  \end{mathpar}

  \medskip
  Subcontext relation for weakening.
  \begin{mathpar}
    (\Ctx[1] \CtxPar \Ctx[3]) \CtxSeq (\Ctx[2] \CtxPar \Ctx[4])
    \SubCtx
    (\Ctx[1] \CtxSeq \Ctx[2]) \CtxPar (\Ctx[3] \CtxSeq \Ctx[4])

    \inferrule{
      \Unr\Ctx
    }{
      \CtxEmpty
      \SubCtx
      \Ctx
    }
  \end{mathpar}

  \medskip
  Laws for context-free session types.
  \begin{mathpar}
    \SSeq\SSkip\S = \S \and
    \SSeq \S\SSkip = \S \and
    \SSeq{\S[1]}{(\SSeq{\S[2]}{\S[3]})} =
    \SSeq{(\SSeq{\S[1]}{\S[2]})}{{\S[3]}} \and
    \SRec\SVar\S = \S{[\SRec\SVar\S/\SVar]} \and
    \SSeq{(\SChoice{\ell:\S[\ell]}{\ell\in L})}\S =
    \SChoice{\ell:\SSeq{\S[\ell]}\S}{\ell\in L} \and
    \SSeq{(\SBranch{\ell:\S[\ell]}{\ell\in L})}\S = \SBranch{\ell:\SSeq{\S[\ell]}\S}{\ell\in L}
  \end{mathpar}

  \medskip
  Laws for duality $\SDual\S$ (there is no dual for $\SRet$ and $\SAcq$)
  \begin{mathpar}
    \SDual\SSkip = \SSkip \and
    \SDual{(\SSeq{\S[1]}{\S[2]})} = \SSeq{\SDual{\S[1]}}{\SDual{\S[2]}}
    \and
    \SDual{\SSend\T} = \SRecv\T \and
    \SDual{\SRecv\T} = \SSend\T \and
    \SDual{(\SChoice{\ell:\S[\ell]}{\ell\in L})} =
    \SBranch{\ell:\SDual{\S[\ell]}}{\ell\in L} \and
    \SDual{(\SBranch{\ell:\S[\ell]}{\ell\in L})} =
    \SChoice{\ell:\SDual{\S[\ell]}}{\ell\in L} \and
    \SDual\STerm = \SWait \and
    \SDual\SWait = \STerm \and
    \SDual{(\SRec\SVar\S)} = \SDual{\S{[\SRec\SVar\S/\SVar]}}
  \end{mathpar}
  \caption{Context and session-type structure}
  \label{fig:syntax-types}
\end{figure}

\Cref{fig:syntax-types} makes context equality sensitive to resource
discipline. Parallel composition is commutative and associative with empty
context as unit. Sequential composition is associative and preserves order.
Unrestricted contexts contract under $\CtxPar$, which is the only source of
implicit duplication. Mobile contexts may move between $\CtxPar$ and
$\CtxSeq$ connectives, reflecting that the type system is not
responsible for ordering access to mobile values. Any unrestricted
context is also mobile (cf.~\cref{fig:bounded-session-types}).

Weakening $\Ctx[1]\SubCtx\Ctx[2]$ is separate from equality. The exchange axiom in the subcontext
relation is the concurrent-Kleene-algebra law that embeds two independent
ordered computations into one ordered computation over parallel
components \cite{DBLP:journals/jlp/HoareMSW11}.
The second subcontext rule admits unused unrestricted assumptions. By
\RuleName{T-Weaken}, a derivation obtained in a smaller context can be used in
any larger context related by $\SubCtx$.

The session-type equations are the standard CFST equations used by type
conversion. Sequential composition has $\SSkip$ as unit and is associative.
Recursive types unfold equi-recursively. Continuations distribute through
internal and external choice, so a suffix can be exposed branchwise before
splitting or comparison.

\begin{figure}
  \begin{mathpar}
    \RuleUUnit \and
    \RuleUArr \and
    \RuleUProd \and
    \RuleUVariant \\
    \RuleUCtxEmpty \and
    \RuleUCtxVar \and
    \RuleUCtxSeq \and
    \RuleUCtxPar
  \end{mathpar}
  \caption{Unrestricted types ($\Unr\T$) and contexts ($\Unr\Ctx$)}
  \label{fig:unrestricted-types}
\end{figure}

\Cref{fig:unrestricted-types} identifies the duplicable fragment. Unit is
unrestricted, products and variants are unrestricted componentwise, and
unrestricted functions have direction $\DirNone$. A context is unrestricted
when all assumptions in it are unrestricted. This predicate is used by
contraction, weakening, and recursive function introduction.

\begin{figure}
  \begin{mathpar}
    \RuleTeSkip \and
    \RuleTeSkipSeq \and
    \RuleTeSkipMu \\
    \RuleTeTerm \and
    \RuleTeWait \and
    \RuleTeRet \and
    \RuleTeSeqOne \\
    \RuleTeRec \and
    \RuleTeBranch \and
    \RuleTeChoice \and
    \RuleTeSeqTwo \\
    \RuleTeBase \and
    \RuleTeArr \and
    \RuleTeAcq \and
    \RuleTeProd \and
    \RuleTeVariant \\
    \RuleTeCtxEmpty \and
    \RuleTeCtxVar \and
    \RuleTeCtxSeq \and
    \RuleTeCtxPar
  \end{mathpar}
  \caption{Mobile types and contexts ($\Mobile\T$, $\Mobile\Ctx$);
    bounded session types ($\Bounded\S$); skipping session types ($\Skips\S$)}
  \label{fig:bounded-session-types}
\end{figure}

\Cref{fig:bounded-session-types} defines the predicates that control
cross-process movement. A skipping session is equivalent to $\SSkip$ and
therefore carries no communication obligation. A bounded session is a finite
obligation: non-recursive bounded sessions terminate in $\STerm$, $\SWait$, or
$\SRet$, while recursive bounded sessions require all finite unfoldings to
reach such a boundary consistently. A mobile session has the form
$\SSeq\SAcq\S$ with $\Bounded\S$; the leading acquire marks a remotely
borrowed endpoint, and boundedness ensures that the borrowed prefix can be
completed without consuming an unbounded residual owned elsewhere.

Mobility extends from sessions to ordinary values. Unit is mobile. Products
and variants are mobile pointwise. Function values are mobile exactly when
their arrow carries the mobile annotation, which in turn means that
the function closes over mobile values. Context mobility is pointwise, so a
process may fork or send only closures whose captured assumptions satisfy the
same discipline.

\begin{figure}
  \begin{mathpar}
    \RuleSNewSkip \and
    \RuleSNewSend \and
    \RuleSNewRecv \and
    \RuleSNewVar \\
    \RuleSNewSeq \and
    \RuleSNewBranch \and
    \RuleSNewChoice \and
    \RuleSNewRec
  \end{mathpar}
  \caption{Session types that are valid for $\CNew$}
  \label{fig:new-session-types}
\end{figure}

The predicate $\SNew\S$ in \cref{fig:new-session-types} restricts the session
types that may be supplied to channel creation. Fresh channels are created for
ordinary communication protocols built from $\SSkip$, sends, receives,
choices, branches, sequencing, and recursion. The synchronization markers
$\SAcq$ and $\SRet$ are not source protocols for $\CNew$; they are introduced
by remote splitting.

\begin{figure}
  \begin{mathpar}
    \RuleCTUnit \and
    \RuleCTNew \and
    \RuleCTFork \and
    \RuleCTLSplit \and
    \RuleCTRSplit \and
    \RuleCTDrop \and
    \RuleCTAcquire \and
    \RuleCTTerm \and
    \RuleCTWait \and
    \RuleCTSend \and
    \RuleCTRecv \and
    \RuleCTSelect \and
    \RuleCTBranch \and
    \RuleCTDiscard
  \end{mathpar}
  \caption{Constant typing ($\HasCType\C\T$)}
  \label{fig:constant-typing}
\end{figure}

\Cref{fig:constant-typing} assign closed types to constants. Channel
creation produces a pair of dual endpoints, each protected by an initial
$\SAcq$ marker and closed respectively by $\STerm$ and $\SWait$. Communication
constants consume exactly the session prefix they implement: send and receive
advance through $\SSend\T$ and $\SRecv\T$, selection advances an internal
choice, and branching returns a variant over external branches. These
constants are impure because they communicate.

The split constants are pure. Local splitting, $\CSplit$, decomposes
$\SSeq{\S[1]}{\S[2]}$ into an ordered pair containing the borrowed prefix and
the residual suffix. The ordering enforces that the prefix must be
consumed before starting with the suffix. Remote splitting, $\CSSplit$, produces an unordered pair
in which the prefix ends in $\SRet$ and the residual begins with $\SAcq$; the
operational semantics uses these markers to synchronize drop and acquire
across processes. The constants $\CDrop$ and $\CAcq$ eliminate the
corresponding markers. Forking expects a mobile impure computation from unit to
unit, and $\CDiscard$ consumes a skipping endpoint.

\begin{figure}
  \begin{mathpar}
    \RuleTConst \and
    \RuleTVar \\
    \RuleTAbsUnr \and
    \RuleTAppUnr \and
    \RuleTAbsLin \and
    \RuleTAppLin \and
    \RuleTAbsLeft \and
    \RuleTAppLeft \and
    \RuleTAbsRight \and
    \RuleTAppRight \and
    \RuleTPairUnord \and
    \RuleTPairOrd \and
    \RuleTLet \and
    \RuleTSeq \and
    \RuleTLetUnord \and
    \RuleTLetOrd \and
    \RuleTInj \and
    \RuleTCase \and
    \RuleTWeaken \and
    \RuleTConv
  \end{mathpar}
  \caption{Expression typing ($\HasType[\Eff]\Ctx\E\T$)}
  \label{fig:typing-declarative}
\end{figure}

The expression typing judgment $\HasType[\Eff]\Ctx\E\T$ in
\cref{fig:typing-declarative} records both resource usage and effects. The
context $\Ctx$ describes the assumptions consumed by $\E$. The effect
annotation is an upper bound: pure expressions may be used where impure
expressions are expected, but not conversely.

Constants and variables are pure and use exactly the resources shown in their
types. Function introduction records how the argument is positioned relative
to the captured context. An unordered function places the argument in parallel
with the captured context; left- and right-ordered functions place it before or
after the captured context; unrestricted recursive functions use direction
$\DirNone$ and may only capture an unrestricted context. Application composes
the contexts of the function and argument according to the arrow direction.
The ordered application rules also enforce the expected purity side condition:
the expression evaluated outside the function's direction must be
pure. In particular, the $\DirLeft$-annotated abstraction puts the
argument left of the captured context. Hence, its application must
evaluate the argument before the function, which implies that the
function must be pure. Any other abstraction can be evaluated
left-to-right as common.

Pair introduction mirrors function application. Unordered pairs combine their
components in parallel. Ordered pairs use a sequential context and require the
second component to be pure. Sequencing requires the subterm to be
unrestricted, so that its result can be ignored. Let and pair elimination expose bindings
to the body using the same ordered or unordered structure as the eliminated
type. Case analysis requires all branches to produce the same result type
under the same surrounding context shape. The final rules close the judgment
under weakening, type conversion, and effect subsumption.

\begin{figure}
  \begin{mathpar}
    \RuleTExpr \and
    \RuleTPar \and
    \RuleTRes \\
    \RuleBindEmp \and
    \RuleBindSeq \and
    \RuleBindSep \and
    \RuleBindVar
  \end{mathpar}
  \caption{Process typing ($\ProcHasType \Ctx {\Proc}$) and binding
    ($\BindGivesCtx {\S} {\BindGroup} {\Ctx}$)}
  \label{fig:process-typing}
\end{figure}

Processes are typed by $\ProcHasType\Ctx\Proc$ in
\cref{fig:process-typing}. An expression process may be an impure computation
of type unit. Parallel composition splits the process context with $\CtxPar$.
Restriction introduces two binder groups at dual session types and adds
the assumptions generated for those binders in parallel to the subprocess.

The auxiliary judgment $\BindGivesCtx\S\BindGroup\Ctx$ relates binder
groups to contexts. A plain sequence of variables follows the sequential
structure of the session type. A separator splits a session
$\SSeq{\S[1]}{\S[2]}$ into a prefix ending in $\SRet$ and a suffix guarded by
$\SAcq$, and places the resulting contexts in parallel. Thus binder groups are
the static representation of the synchronization boundary introduced by
remote borrowing. This judgment is nondeterministic as session types
are subject to conversion.

%%% Local Variables:
%%% mode: latex
%%% TeX-master: "../popl2027.tex"
%%% End:
